\documentclass{article}
\usepackage[UTF8]{ctex}
\usepackage{amsmath,amssymb,amsfonts}
\usepackage{enumitem}
\usepackage{geometry}
\geometry{a4paper, margin=1in}
% === 错题本排版优化（不参与题目标识） ===
\setlength{\parindent}{0pt}
\setlength{\parskip}{0.35em plus 0.10em minus 0.08em}
\setlength{\emergencystretch}{3em}
\setlist{itemsep=0.25em, topsep=0.40em, parsep=0pt}
\allowdisplaybreaks[2]
% === 错题本排版优化结束 ===
\begin{document}
\section*{华南理工大学2024年数学分析真题}
\section*{155．华南理工大学2024年数学分析试卷}

\subsection{题目1}
1．（13 分）已知 \(a, b>0\) ，且 \(c=\lim _{x \rightarrow 0}\left(\frac{4}{\ln (1+x)}-\frac{4}{x}\right)\) ，定义函数

\[
f(x)= \begin{cases}\frac{\sin (a x)}{x}, & x<0 ; \\ c, & x=0 ; \\ (1+b x)^{\frac{1}{x}}, & x>0 .\end{cases}
\]

若 \(f(x)\) 在 \((-\infty,+\infty)\) 内处处连续，求常数 \(a, b\) 和 \(c\) 。

\subsection{题目2}
2．（13 分）设 \(f(x)\) 在 \([a, b]\) 上连续，在 \((a, b)\) 内可导．证明：存在 \(\xi \in(a, b)\) ，使得

\[
f(\xi)+f^{\prime}(\xi)=\frac{e^{b} f(b)-e^{a} f(a)}{e^{b}-e^{a}}
\]

\subsection{题目3}
3．（13 分）证明函数

\[
f(x, y)= \begin{cases}x y \cos \frac{1}{\sqrt{x^{2}+y^{2}}}, & (x, y) \neq(0,0) \\ 0, & (x, y)=(0,0)\end{cases}
\]

在点 \((0,0)\) 连续且偏导数存在，但是偏导数在点 \((0,0)\) 不连续，而 \(f(x, y)\) 在点 \((0,0)\) 可微．

\subsection{题目4}
4．（13 分）将方程 \(\frac{\partial^{2} z}{\partial x^{2}}+\frac{\partial^{2} z}{\partial y^{2}}=0\) 变换为极坐标形式。

\subsection{题目5}
5．（14 分）计算曲线积分 \(\int_{L} y \mathrm{~d} x+z \mathrm{~d} y+x \mathrm{~d} z\) ，其中 \(L\) 是曲线

\[
\frac{x^{2}}{a^{2}}+\frac{y^{2}}{b^{2}}+\frac{z^{2}}{c^{2}}=1, \frac{x}{a}+\frac{z}{c}=1, x \geq 0, y \geq 0, z \geq 0 .
\]

（ \(a>0, b>0, c>0\) 为常数）从点 \((a, 0,0)\) 到 \((0,0, c)\) ．

\subsection{题目6}
6．（14 分）证明黎曼 zeta 函数 \(\zeta(x)=\sum_{n=1}^{\infty} \frac{1}{n^{x}}\) 在 \((1,+\infty)\) 上连续且无穷次可微．

\subsection{题目7}
7．（14 分）解答如下问题：

\begin{itemize}
  \item[（1）] 用定义证明 \(\lim _{n \rightarrow \infty} \sin n\) 不存在。
  \item[（2）] 证明数列 \(\left\{x_{n}\right\}\) 极限存在，其中
\[
x_{n}=\left(1-\frac{1}{2^{2}}\right)\left(1-\frac{1}{3^{2}}\right) \cdots\left(1-\frac{1}{(n+1)^{2}}\right)=\prod_{k=1}^{n}\left(1-\frac{1}{(k+1)^{2}}\right), n=1,2, \cdots .
\]
并求极限 \(\lim _{n \rightarrow \infty} x_{n}\) ．
\end{itemize}

\subsection{题目8}
8．（14 分）设函数 \(f(x)\) 在 \(\left(x_{0}-r, x_{0}+r\right)\) 上无穷次可导（其中 \(r \in(0,+\infty)\) ），且存在常数 \(M>0\) ，使得

\[
\left|f^{(n)}(x)\right| \leq M \frac{n!}{r^{n}}, \forall x \in\left(x_{0}-r, x_{0}+r\right), \forall n=0,1,2, \cdots .
\]

\begin{itemize}
  \item[（1）] 证明 \(f(x)\) 在 \(x_{0}\) 点的 Taylor 级数 \(\sum_{n=0}^{\infty} \frac{f^{(n)}\left(x_{0}\right)}{n!}\left(x-x_{0}\right)^{n}\) 在 \(\left(x_{0}-r, x_{0}+r\right)\) 上逐点收敛于 \(f(x)\) 。
  \item[（2）] 证明 \(\sum_{n=0}^{\infty} \frac{f^{(n)}\left(x_{0}\right)}{n!}\left(x-x_{0}\right)^{n}\) 在 \(\left(x_{0}-r, x_{0}+r\right)\) 上内闭一致收敛于 \(f(x)\) ．
\end{itemize}

\subsection{题目9}
9．（14 分）设 \(f(x)\) 在 \([a, b]\) 上黎曼可积，证明黎曼引理

\[
\lim _{\lambda \rightarrow+\infty} \int_{a}^{b} f(x) \cos (\lambda x) \mathrm{d} x=0
\]

\subsection{题目10}
10．（14 分）证明 \(\int_{0}^{+\infty} \frac{\sin x y}{x} \mathrm{~d} x\) 关于 \(y\) 在 \(\left[y_{0},+\infty\right)\left(y_{0}>0\right)\) 上一致收敛，但在 \((0,+\infty)\) 上非一致收敛。

\subsection{题目11}
11．（14 分）设 \(f(x)\) 在 \([a, b]\) 上分段连续，即存在 \([a, b]\) 的一个有限分割 \(a=x_{0}<x_{1}<x_{2}<\cdots<x_{n}=\) \(b\)（其中 \(n\) 为固定整数），使得 \(f\) 在每个区间 \(\left(x_{i-1}, x_{i}\right)\) 上连续且分点 \(x_{i}\) 处都存在左右极限．证明 \(f(x)\)在 \([a, b]\) 上黎曼可积．
\end{document}
